\color{unimain}{\section{Описание}}
\color{black}

\color{uniblue}{\subsection{Назначение}}
\color{black}

\par
\sloppy

Устройство серии ЮНИТ-М300 применяется в схемах вторичной коммутации распределительных устройств 6-35 кВ с переменным, постоянным и выпрямленным оперативным током. Устройство предназначено для осуществления функций релейной защиты и автоматики первичного оборудования, измерения, регистрации, осциллографирования и сигнализации на объектах энергетики классов напряжения 6–35 кВ. Устройство выпускается в нескольких типоисполнениях по набору функций защит и автоматики. Подробнее о функционале устройства конкретного типоисполнения описано в РЭ на него. 
Существует возможность выполнения функций дистанционного управления и сбора данных при подключении устройства к информационной сети энергообъекта.
Взаимодействие с устройством может быть реализовано при помощи подключаемого устройства ЮНИТ-ИЧМ или через ПК с установленным ПО «ЮНИТ-Сервис». Руководство по устройству ЮНИТ-ИЧМ представлено в документе RU.37182817.00001.01.34.09 – «Устройство «ЮНИТ-ИЧМ». Руководство по эксплуатации». Руководство по использованию ПО представлено в документе RU.37182817.00001.01.34.09 – «Программное обеспечение мониторинга и параметрирования «ЮНИТ-Сервис». Руководство пользователя». 
Устройство устанавливается в релейных отсеках ячеек КРУ, КРУН, камер КСО, на панелях и в шкафах релейных залов электростанций и подстанций.
Устройства серии ЮНИТ-М300 построены на микропроцессорной элементной базе.

\subsection{Технические характеристики устройства}
\subsubsection{Технические данные}
Технические данные устройства указаны в таблицах 1.1 – 1.6.

Таблица 1.1 - Модуль источника питания P01